
\clearpage
\VerdierPhysicalPageAnchor{0124}%
\begingroup
\small
\setlength{\emergencystretch}{1.5em}%
\setlength{\abovedisplayskip}{4pt plus 2pt minus 2pt}%
\setlength{\belowdisplayskip}{4pt plus 2pt minus 2pt}%

\VerdierSemanticAnchor{ch2-1-3-2-triangle-morphism-conditions}
\noindent が存在し、$(\operatorname{id}_{X_1},u_1,m'_1)$ と
$(u_3,\operatorname{id}_{X_3},m'_3)$ は三角の射であり、さらに三角

\VerdierSemanticAnchor{ch2-1-3-2-connecting-triangle}
\[
  Z'_3\xrightarrow{m'_1}Z'_2\xrightarrow{m'_3}Z'_1
  \xrightarrow{\,v'_3[1]w'_1\,}Z'_3[1]
\]
\noindent は完全三角である。

\VerdierSemanticAnchor{ch2-1-3-2-comparison-isomorphisms}
さらに (1.2.4) と (1.2.14) により、
$(\operatorname{id},\operatorname{id},i_j)$ が三角の同型となるような同型
$i_j\colon Z'_j\longrightarrow Z_j$（$1\leq j\leq 3$）が存在する。そこで

\VerdierSemanticAnchor{ch2-1-3-2-transported-morphisms}
\[
  m_1=i_2m'_1i_3^{-1}\colon Z_3\longrightarrow Z_2,
  \qquad
  m_3=i_1m'_3i_2^{-1}\colon Z_2\longrightarrow Z_1.
\]

\VerdierSemanticAnchor{ch2-1-3-2-proof-complete}
とおけばよい。射 $m_1$ と $m_3$ は公理 TRIV が要求する性質をもつ。

\VerdierSemanticAnchor{ch2-1-3-3}
\medskip
\noindent\textbf{1.3.3.}\quad
$C$ を\emph{加法}圏とする。$C$ の各対象 $X$ に、次数零の成分だけが $X$
であり、他の成分がすべて零である $C$ の複体を対応させる。これにより $C$
から $\operatorname{comp}(C)$ への関手が定まり、さらに標準関手
$\operatorname{comp}(C)\longrightarrow K(C)$ と合成して、$C$ から $K(C)$
への関手が定まる。

\VerdierSemanticAnchor{ch2-1-3-3-full-embeddings}
これらの関手は充満忠実で、対象上単射である。したがって $C$ は
$\operatorname{comp}(C)$ と $K(C)$ の充満部分圏として実現される。以下では、
常にこれらの関手によって圏 $C$ を $\operatorname{comp}(C)$ と $K(C)$ の
充満部分圏とみなす。

\VerdierSemanticAnchor{ch2-1-3-4}
\medskip
\noindent\textbf{1.3.4.}\quad
$C$ を\emph{アーベル}圏とする。$C$ の各複体 $X$ に対し、$C$ の対象

\VerdierSemanticAnchor{ch2-1-3-4-1}
\[
  H^0(X)=\operatorname{Ker}(d_X^0)/\operatorname{Im}(d_X^{-1}).
  \tag{1.3.4.1}
\]
\noindent を $H^0(X)$ と書く。

\VerdierSemanticAnchor{ch2-1-3-4-induced-morphism}
複体の各射 $X\xrightarrow{f}Y$ に対し、
\[
  H^0(f)\colon H^0(X)\longrightarrow H^0(Y)
\]
\noindent を $f$ が誘導する射とする。これにより関手

\VerdierSemanticAnchor{ch2-1-3-4-comp-functor}
\[
  H^0\colon\operatorname{comp}(C)\longrightarrow C
\]
\noindent が定まり、これを\emph{第零コホモロジー対象}と呼ぶ。この関手は
ホモトピックな射を等しい射に送るので、同じく

\VerdierSemanticAnchor{ch2-1-3-4-k-functor}
\[
  H^0\colon K(C)\longrightarrow C.
\]
\endgroup
